\section{Literature Review}

This chapter reviews the literature relevant to PyTrade's architecture. The proposed system separates temporal signal extraction for individual assets from cross-sectional portfolio allocation. The motivation for this separation comes from several established lines of research. Reinforcement learning provides a framework for sequential portfolio decisions. Recurrent models provide memory for partially observed time series. Attention provides a mechanism for relating multiple assets. Hierarchical and modular reinforcement learning provides principles for combining specialized decision components. The review also emphasizes methodological risks in financial machine learning, because a more complex architecture does not by itself produce more reliable evidence.

\subsection{Quantitative Finance and Deep Reinforcement Learning}

Quantitative finance has traditionally used statistical models, optimization, and rule-based trading systems to describe prices, estimate risk, and construct portfolios. These methods remain important because they provide explicit assumptions and useful baselines. Their limitations become more visible when the decision is sequential. A trading system must decide not only which assets to hold, but also when to trade, how to size multiple bets, and how current positions affect future choices. Transaction costs, cash constraints, and changing risk exposures make this a control problem rather than a sequence of independent forecasts.

Reinforcement learning formulates such a problem through states, actions, transitions, and rewards. The agent observes a state, selects an action, receives a reward, and observes a new state. Portfolio management is consequently often modeled as a continuous-control problem in which actions represent positions or portfolio weights. The objective can include return, risk, and trading costs. A portfolio-management study by Wu et al. illustrates this use of a reinforcement-learning framework for portfolio decisions \cite{Junfeng.2024}. The FinRL work further shows the importance of an explicit financial environment, data pipeline, and evaluation procedure when deep reinforcement learning is applied to automated trading \cite{liu2022finrldeepreinforcementlearning}.

The action representation determines what the policy must learn. Discrete-action systems can express decisions such as buy, sell, or hold. They are comparatively simple to interpret, but they require a discretization of trade size and may become unwieldy when several assets and position sizes are considered simultaneously. Continuous-action systems can represent position changes or portfolio weights directly. They fit the allocation problem more naturally, but they also create a larger action space and make constraint handling part of the environment design. \textit{[Literature needed: a source that directly compares discrete-action deep reinforcement learning with continuous-action policy-gradient methods for portfolio management. No specific candidate is currently available in the bibliography.]}

Proximal Policy Optimization (PPO) is a policy-gradient method designed to limit excessively large policy updates. Its clipped objective constrains the policy ratio in the main update and thereby reduces the risk of destructive changes during training \cite{JohnSchulmanFilipWolskiPrafullaDhariwalAlecRadfordOlegKlimov.2017}. This property motivates the use of PPO for the continuous actions in PyTrade. PPO does not, however, make financial learning stable automatically. It does not remove non-stationarity from prices, enforce portfolio constraints, or prevent the policy from exploiting an unrealistic simulator. These responsibilities remain with the state representation, environment, reward, and evaluation protocol.

The distinction between a trading strategy and a meta-strategy is useful in this context. A trading strategy produces decisions from market information. A meta-strategy determines how strategies are selected, combined, sized, or switched. In a modular portfolio system, the SAA can be viewed as a local strategy that produces an asset-level proposal, while the PAA acts as a meta-strategy that decides how those proposals should be combined under portfolio-wide constraints. \textit{[Literature needed: a finance source defining and evaluating meta-strategies for combining trading strategies or signals. No specific candidate is currently available in the bibliography.]}

López de Prado describes three broad sources of error in machine-learning models: bias, variance, and noise \cite{LopezdePrado.2018}. Bias arises when the model is too restrictive to represent the relevant relationship. Variance arises when the model adapts too closely to the particular sample used for development. Noise is irreducible variation or measurement error that cannot be learned reliably. Increasing the complexity of a hierarchical model may reduce bias by giving different components specialized functions. It may also increase variance because more parameters, training choices, and interfaces can be tuned. The decomposition is therefore a hypothesis about how to manage error, not an argument that a larger architecture must perform better.

The same methodological perspective motivates safeguards against the seven sins of quantitative investing discussed by López de Prado: survivorship bias, look-ahead bias, storytelling, data mining and statistical snooping, ignoring transaction costs, mishandling outliers, and unjustified shorting assumptions \cite{LopezdePrado.2018}. These risks map directly onto PyTrade. The asset universe and historical data must not exclude failed or delisted assets without explanation. Features and rewards must use only information available at the decision time. A plausible narrative about a market regime is not evidence of a tradable signal. Repeated experiments and hyperparameter searches require validation controls. Commissions, spreads, and market impact must be represented. Outliers require robust analysis rather than automatic removal. Finally, short positions require explicit inventory, financing, and risk assumptions.

Financial reinforcement-learning results are particularly sensitive to the simulator and evaluation design. An agent can exploit an unrealistic execution rule, leakage in feature construction, or a favorable market period. A meaningful evaluation must therefore separate training and unseen evaluation periods, model trading frictions, and compare the learned policy with cash, buy-and-hold, and other simple baselines \cite{LopezdePrado.2018}. These requirements motivate the experimental design of the custom environment PyTrade. The purpose is not to claim that reinforcement learning discovers universal alpha. It is to test whether a modular policy can learn temporal signals and portfolio-level allocation with fewer interfering responsibilities.

\subsection{Temporal Memory in Financial Series}


Market observations are not necessarily Markovian when only a finite set of current features is provided. Relevant information may be distributed across several previous observations. Recurrent neural networks address this issue by maintaining a hidden state that summarizes prior inputs. LSTMs add gated memory and are designed to retain or discard information over longer sequences. Recurrent value-based learning has explicitly been studied for partially observable decision processes, where memory is needed to infer hidden state information \cite{DBLP:journals/corr/HausknechtS15}.

For financial series, temporal memory has two roles. First, it can filter short-term noise and combine observations into a representation of recent behavior. Second, it can help identify regimes whose evidence is spread over time. These roles are related but not identical. A representation that is useful for detecting a local price pattern need not be the best representation for deciding how capital should be distributed across assets. The two-level design in PyTrade assigns local temporal processing to the SAA and leaves cross-asset comparison to the PAA.

Stationarity creates a second difficulty. Raw prices are usually non-stationary, while simple returns are closer to stationary but remove much of the level information. Integer differencing can make a series more suitable for statistical modeling, but it can also discard useful long-range dependence. Financial machine-learning methodology therefore discusses fractional differentiation as a compromise between stationarity and memory preservation \cite{LopezdePrado.2018}. This idea is relevant to PyTrade's feature engineering: the input should contain information that is sufficiently stable for learning without erasing economically meaningful temporal structure.

Feature engineering does not guarantee predictive power. It changes the information available to the policy and can introduce leakage if it uses future observations or incorrectly aligned data. In a reinforcement-learning environment, features also interact with the state transition and reward definition. López de Prado connects feature design to the wider problem of controlling model error. A feature can reduce bias by exposing a relevant relationship, but a large collection of candidate features can increase variance through data mining and selection after repeated testing \cite{LopezdePrado.2018}. Look-ahead bias can enter through incorrectly aligned rolling statistics. Outliers can dominate transformations. Storytelling can encourage the selection of features because they fit an attractive explanation after the fact.

The Temporal Fusion Transformer combines attention with mechanisms for temporal selection and interpretable multi-horizon forecasting \cite{lim2020tft}. Its forecasting objective differs from PyTrade's control problem, but it supports the broader idea that temporal inputs can be filtered and combined selectively rather than passed to a model as an undifferentiated vector. \textit{[Literature needed: a controlled financial comparison of RNNs, LSTMs, and temporal convolutional networks. No specific candidate is currently available in the bibliography.]}

The SAA uses recurrent memory to produce a compact single-asset representation. That representation can include a proposed action and hidden-state information for the portfolio-level module. This creates a meaningful interface between the two levels, but it also creates a distribution-shift risk. During independent training, an SAA may observe a particular local portfolio state. During deployment, several SAAs interact through shared cash and portfolio constraints. Shadow portfolios are therefore used in PyTrade to keep the recurrent local state from being driven directly by the changing actions of other assets.

\subsection{Attention Mechanisms}


Attention assigns input-dependent importance to elements of a sequence or set. The Transformer architecture showed how self-attention can model relationships between tokens without requiring recurrence as the main sequence mechanism \cite{vaswani2023attentionneed}. Each token can contribute information to the representation of another token. In a time series, tokens may represent time steps. In a portfolio, they can represent assets. This distinction is important for PyTrade: the PAA uses attention primarily to model cross-sectional relationships, while the SAA uses recurrence to model local temporal context.

Self-attention is useful for allocation because an asset cannot be evaluated independently of the alternatives available in the same portfolio. A proposed position in an equity index may have a different meaning when another equity index, gold, or oil already receives a large allocation. Attention can condition the representation of one asset on the other asset tokens and on a portfolio token containing shared state. This allows the allocator to represent substitution, concentration, and cross-sectional context without requiring a separate hand-written interaction rule for every pair of assets.

Attention also has limitations. Its flexibility can increase variance and make the learned allocation difficult to interpret. With limited financial data, a high-capacity attention model may fit sample-specific correlations instead of stable relationships. Attention weights should not automatically be interpreted as causal importance or reliable economic explanations. The model still requires out-of-sample evaluation, friction modeling, and controls for repeated experimentation. The attention mechanism is a representation tool, not a guarantee of diversification or risk control.

Attention has also entered reinforcement learning. Hu et al. review the development of Transformer-based reinforcement learning and describe how attention-based sequence models can represent trajectories, histories, and decision context \cite{Hu.2024}. Their review supports the use of attention for conditioning decisions on multiple pieces of context, while also showing that Transformer-based reinforcement learning is a broad family of approaches rather than one single algorithm. Kachaev et al. discuss Transformers in online reinforcement learning for continuous control, which is relevant to the continuous action setting of portfolio allocation \cite{NikitaKachaev.2025}. The transfer to finance remains non-trivial because online control benchmarks generally do not include market non-stationarity, trading frictions, or shared portfolio constraints.

In quantitative finance, attention-based models have been proposed for dynamic portfolio optimization. Xue and Ye study attention-enhanced reinforcement learning for this setting and explicitly discuss transaction costs and risk management in the context of dynamic portfolio optimization \cite{Xue.2026}. This is close to the portfolio-level component of PyTrade. However, an attention head that processes multiple assets does not by itself provide a dedicated temporal expert for each asset. The temporal and cross-sectional functions may still be learned jointly, with the associated risk that one function dominates or that the representation becomes difficult to interpret.

Attention-based portfolio research is closer to the PAA setting. Xue and Ye study attention-enhanced reinforcement learning for dynamic portfolio optimization and explicitly include transaction costs and risk management in the problem formulation \cite{Xue.2026}. Coulter's work is also directly framed around neural portfolio allocators, cross-asset optimization, attention-based Transformers, and multi-agent methods \cite{NathanielCoulter.2025}. \textit{[Please verify the publication status and exact empirical claims of the Coulter entry before relying on it in the final thesis, because its current BibTeX record does not specify a publication venue.]}

The proposed PAA represents the portfolio as a portfolio token and a set of asset tokens. An asset token can contain current market features, local shadow-portfolio information, an SAA action, and a recurrent latent state. The portfolio token provides shared context. Attention then coordinates the local proposals into target portfolio weights. This structure gives the allocator access to the problem, the proposed local solution, and information about the local temporal context. It also creates a feature-dilution risk: the raw features, actions, hidden states, and portfolio variables may compete for a limited representation. Ablation experiments that remove the latent state or the SAA action are therefore needed to identify which information contributes to performance.

\subsection{Hierarchical \& Modular RL}


Hierarchical reinforcement learning decomposes a complex decision process into multiple levels or policies. A higher-level component can select goals, sub-tasks, or allocations, while lower-level components execute specialized behavior. Modular reinforcement learning uses a related idea: different modules learn reusable functions and are combined by a larger policy. The potential benefit is reduced interference. A module that extracts a temporal signal for one asset does not need to learn the full portfolio-allocation problem at the same time.

This decomposition also introduces new problems. The interfaces between modules must have a clear meaning and compatible time scales. A higher-level policy can become dependent on artifacts of a lower-level policy. Freezing or copying a module changes the distribution of the inputs seen by the downstream policy. In addition, independently trained modules may produce actions that are locally sensible but jointly infeasible. These issues are especially important in trading because cash, positions, and transaction costs couple all assets through the shared portfolio.

The hybrid LSTM and PPO architecture is relevant because it combines recurrent temporal processing with a policy-optimization method for dynamic portfolio decisions \cite{JunKevin.2025}. It motivates studying recurrence and PPO together, but the roles of the two components should remain conceptually distinct. An LSTM can provide memory, while PPO determines how the policy is updated; neither component alone defines a hierarchical portfolio architecture. Regime-aware approaches provide another motivation for specialization. Raj's work explicitly considers adaptive and regime-aware reinforcement learning for portfolio optimization \cite{GabrielNixonRaj.2025}. A regime-aware policy may adapt its allocation behavior when market conditions change, but regime information must be estimated from historical observations and can be uncertain. \textit{[Please verify the exact methodology and results of the Raj preprint before making stronger claims about regime detection or adaptation. The current BibTeX entry identifies it as an arXiv paper but does not provide peer-review information.]}

The main strength of modularity is not simply that it adds more neural components. It creates an experimental distinction between temporal extraction and cross-sectional allocation. A raw-feature allocator can be compared with an allocator augmented by frozen SAA outputs. The SAA action can be replaced by structured noise to test whether the PAA benefits from historical information or merely from additional dimensions. A monolithic recurrent allocator can test whether joint end-to-end memory performs better than the modular interface. These comparisons turn the architectural proposal into an empirical test rather than an assumption about superiority.

PyTrade implements the hierarchy through a Single-Asset Agent and a Portfolio Allocator Agent. The SAA learns temporal behavior for one instrument while training across a rotating asset set. The PAA receives the resulting proposals and coordinates them across assets. Shadow portfolios preserve the local state of each SAA and reduce direct feedback from the actions of other assets into its recurrent memory. The hierarchy must still be evaluated against its costs: additional training stages, frozen representations, possible information loss at the interface, and more opportunities for variance and data snooping. Under a controlled walk-forward protocol, the central question is whether this separation improves stability and efficiency compared with simpler alternatives.